% ============================================================
% report.tex — 재생에너지 입찰제 도입에 따른 SMP 변화 분석
% 한국 육지계통 개별 발전기 122기 경제급전 모형
% ============================================================
% 컴파일: xelatex report.tex (kotex 사용 시)
%         또는 pdflatex report.tex (kotex 대신 CJK 사용 시)
% ============================================================
\documentclass[11pt, a4paper]{article}

% ── 한국어 지원 ──
\usepackage{kotex}              % XeLaTeX/LuaLaTeX용 한국어 지원
% kotex이 동작하지 않는 환경에서는 아래를 사용:
% \usepackage[utf8]{inputenc}
% \usepackage{CJKutf8}

% ── 수학 패키지 ──
\usepackage{amsmath}
\usepackage{amssymb}
\usepackage{mathtools}
\usepackage{bm}

% ── 레이아웃 ──
\usepackage[margin=2.5cm]{geometry}
\usepackage{setspace}
\onehalfspacing

% ── 표·그림 ──
\usepackage{booktabs}
\usepackage{array}
\usepackage{multirow}
\usepackage{longtable}
\usepackage{tabularx}

% ── 기타 ──
\usepackage{graphicx}
\usepackage{float}
\usepackage{enumitem}
\usepackage{hyperref}
\usepackage{cleveref}
\usepackage{xcolor}
\usepackage{algorithm}
\usepackage{algorithmic}
\usepackage{listings}
\usepackage{caption}
\usepackage{subcaption}
\usepackage{fancyhdr}
\usepackage{titlesec}

% ── 하이퍼링크 설정 ──
\hypersetup{
    colorlinks=true,
    linkcolor=blue!60!black,
    citecolor=green!50!black,
    urlcolor=blue!70!black,
}

% ── 코드 스타일 ──
\lstset{
    language=Julia,
    basicstyle=\ttfamily\small,
    keywordstyle=\color{blue!70!black}\bfseries,
    commentstyle=\color{green!50!black}\itshape,
    stringstyle=\color{red!60!black},
    numbers=left,
    numberstyle=\tiny\color{gray},
    frame=single,
    breaklines=true,
    captionpos=b,
}

% ── 헤더/푸터 ──
\pagestyle{fancy}
\fancyhf{}
\fancyhead[L]{\small 재생에너지 입찰제 SMP 분석}
\fancyhead[R]{\small \thepage}
\fancyfoot[C]{\small PSE Project -- 개별 발전기 122기 모형}

% ── 수학 연산자 정의 ──
\DeclareMathOperator*{\argmin}{argmin}
\DeclareMathOperator*{\argmax}{argmax}
\newcommand{\mw}{\text{MW}}
\newcommand{\mwh}{\text{MWh}}
\newcommand{\won}{\text{원}}
\newcommand{\kwh}{\text{kWh}}
\newcommand{\gcal}{\text{Gcal}}

% ============================================================
\title{
    \vspace{-2cm}
    \textbf{재생에너지 입찰제 도입에 따른 \\
    한국 육지계통 SMP 변화 분석} \\[0.5cm]
    \large 개별 발전기 122기 경제급전(ED) 모형 \\
    수학적 정식화 및 구현 보고서
}
\author{
    전력시스템 경제학 프로젝트
}
\date{\today}

\begin{document}
\maketitle
\thispagestyle{empty}

\begin{abstract}
본 보고서는 한국 육지(mainland) 전력계통에서 재생에너지 입찰제 도입 시
계통한계가격(SMP)에 미치는 영향을 분석하기 위해 개발된 경제급전(Economic Dispatch, ED)
모형의 수학적 정식화를 문서화한다.
모형은 MATPOWER KPG193 기반 개별 발전기 122기를 포함하며,
Basic ED, Pre-revision ED, Post-revision ED의 3단계로 구성된다.
각 단계의 목적함수, 제약조건, SMP 도출 방법을 엄밀히 기술하고,
Price Adder 보정 알고리즘, 구간별 선형 비용 근사, 재생에너지 6블록 입찰 구조,
Beta 혼합 몬테카를로 시뮬레이션 등의 수학적 세부사항을 제시한다.
구현은 Julia/JuMP 프레임워크와 HiGHS 선형계획 솔버를 사용하였으며,
2024년 8,784시간 실데이터를 기반으로 검증하였다.
\end{abstract}

\tableofcontents
\newpage

% ============================================================
% 1. 서론
% ============================================================
\section{서론}\label{sec:intro}

\subsection{프로젝트 개요}

본 프로젝트는 한국 전력시장에서 재생에너지(태양광, 풍력) 발전사업자가
직접 입찰에 참여하는 제도 변경이 계통한계가격(System Marginal Price, SMP)에
미치는 영향을 정량적으로 분석하는 것을 목적으로 한다.

현행 한국 전력시장(CBP: Cost-Based Pool)에서 재생에너지는
음의 부하(negative load)로 처리되어 순수요(net demand)에서 차감된다.
그러나 재생에너지 비중 증가에 따라 출력제한(curtailment)이 빈번해지고,
시장가격 형성에 재생에너지의 한계비용이 반영되지 못하는 문제가 발생한다.
이에 따라 재생에너지 입찰참여 제도가 검토되고 있으며,
본 모형은 이러한 제도 변경의 SMP 효과를 사전 평가한다.

\subsection{한국 육지계통 개요}

분석 대상은 한국 육지(mainland) 전력계통으로, 다음과 같은 특성을 갖는다:
\begin{itemize}[noitemsep]
    \item \textbf{발전기}: MATPOWER KPG193 기반 개별 발전기 122기
    \item \textbf{연료원}: 원자력(Nuclear), 석탄(Coal), LNG, 유류(Oil) 등
    \item \textbf{재생에너지}: 태양광(PV) 및 풍력(Wind) --- 2024년 설비용량 기준
    \item \textbf{시장제도}: 비용 기반 풀(CBP) --- 열소비율이 gencost에 내포
    \item \textbf{분석기간}: 2024년 1월--12월 (8,784시간, 윤년)
\end{itemize}

\subsection{모형 구조 개요}

모형은 3단계 경제급전 모형으로 구성된다:
\begin{enumerate}[noitemsep]
    \item \textbf{Basic ED}: 단순 선형비용 경제급전 (벤치마크)
    \item \textbf{Pre-revision ED}: 현행 시장규칙 반영 (최소출력, 램프, must-run, Price Adder)
    \item \textbf{Post-revision ED}: 재생에너지 입찰 블록 포함 (6블록, BidFloor, RE Pmin)
\end{enumerate}
각 단계의 SMP는 수급균형 제약의 쌍대변수(dual variable)로 도출되며,
$\Delta\text{SMP} = \text{SMP}_{\text{post}} - \text{SMP}_{\text{pre}}$를 핵심 분석 지표로 사용한다.

\subsection{표기법}

본 보고서에서 사용하는 주요 표기법은 \Cref{tab:notation}에 정리한다.

\begin{table}[H]
\centering
\caption{주요 표기법}\label{tab:notation}
\small
\begin{tabularx}{\textwidth}{>{\raggedright}p{3cm} X}
\toprule
\textbf{기호} & \textbf{설명} \\
\midrule
$\mathcal{G} = \{1, \ldots, G\}$ & 열발전기 집합 ($G=122$) \\
$\mathcal{T} = \{1, \ldots, T\}$ & 시간 인덱스 집합 ($T=24$, 1일 단위) \\
$\mathcal{K} = \{1, \ldots, K\}$ & 재생에너지 입찰 블록 집합 ($K=6$) \\
$\mathcal{S} = \{1, \ldots, S\}$ & 구간별 선형 근사 세그먼트 집합 ($S=4$) \\
$p_{g,t}$ & 발전기 $g$의 시간 $t$ 발전량 [MW] \\
$r_{k,t}$ & 입찰 블록 $k$의 시간 $t$ 낙찰량 [MW] \\
$D_t$ & 시간 $t$의 총 수요 [MW] \\
$\text{RE}_t$ & 시간 $t$의 재생에너지 발전량 [MW] \\
$\lambda_t$ & 시간 $t$의 SMP [원/MWh] (수급균형 쌍대변수) \\
$P_g^{\min}, P_g^{\max}$ & 발전기 $g$의 최소/최대 출력 [MW] \\
$\text{RU}_g, \text{RD}_g$ & 발전기 $g$의 상향/하향 램프 한계 [MW/h] \\
$A_{g,t}$ & 발전기 $g$의 시간 $t$ Price Adder [원/MWh] \\
$\bar{R}_{k,t}$ & 입찰 블록 $k$의 시간 $t$ 공급가능량 [MW] \\
$b_{k,t}$ & 입찰 블록 $k$의 시간 $t$ 입찰가격 [원/MWh] \\
\bottomrule
\end{tabularx}
\end{table}


% ============================================================
% 2. 데이터 구조
% ============================================================
\section{데이터 구조}\label{sec:data}

\subsection{발전기 데이터}

개별 발전기 122기의 데이터는 \texttt{generators.csv}에서 로딩되며,
각 발전기는 다음 속성을 갖는 \texttt{ThermalGenerator} 구조체로 표현된다:

\begin{table}[H]
\centering
\caption{ThermalGenerator 구조체 필드}\label{tab:thermal_gen}
\small
\begin{tabular}{lll}
\toprule
\textbf{필드} & \textbf{단위} & \textbf{설명} \\
\midrule
\texttt{name}           & ---   & 발전기 식별자 (예: ``LNG\_001'', ``Nuclear\_003'') \\
\texttt{fuel}           & ---   & 연료원 (``Nuclear'', ``Coal'', ``LNG'', ``Oil'') \\
\texttt{pmin}           & MW    & 최소출력 \\
\texttt{pmax}           & MW    & 최대출력 \\
\texttt{ramp\_up}       & MW/h  & 상향 램프 한계 \\
\texttt{ramp\_down}     & MW/h  & 하향 램프 한계 \\
\texttt{must\_run}      & bool  & must-run 여부 (원자력 = true) \\
\texttt{marginal\_cost} & 원/MWh & 중점 한계비용 (Basic ED용) \\
\bottomrule
\end{tabular}
\end{table}

한국 CBP 시장의 특성상 열소비율(Heat Rate)은 gencost의 2차 비용함수에 내포되어 있으며,
변동운영비(VOM)는 별도 정산되어 발전비용에 포함되지 않는다.

\subsection{Gencost: 2차 비용함수}\label{sec:gencost}

각 발전기의 비용함수는 \texttt{gencost.csv}에서 로딩되며, 다음과 같은 2차 함수 형태를 갖는다:
\begin{equation}\label{eq:gencost}
    C_g(P) = a_g P^2 + b_g P + c_g \qquad [\text{천원/h}]
\end{equation}
여기서:
\begin{itemize}[noitemsep]
    \item $a_g$: 2차 계수 [천원/MW$^2$h]
    \item $b_g$: 1차 계수 [천원/MWh]
    \item $c_g$: 상수항 (무부하비) [천원/h]
    \item $P$: 발전출력 [MW]
\end{itemize}
gencost 계수는 딕셔너리 형태 $\texttt{gencost} : \text{발전기명} \to (a_g, b_g, c_g)$로 저장된다.

\subsection{한계비용(MC) 도출}\label{sec:mc_derivation}

비용함수 \eqref{eq:gencost}로부터 한계비용(Marginal Cost)을 미분으로 도출한다:
\begin{equation}\label{eq:mc}
    \text{MC}_g(P) = \frac{dC_g}{dP} = 2a_g P + b_g \qquad [\text{천원/MWh}]
\end{equation}
단위 변환(천원 $\to$ 원):
\begin{equation}\label{eq:mc_won}
    \boxed{\text{MC}_g(P) = (2a_g P + b_g) \times 1000 \qquad [\text{원/MWh}]}
\end{equation}

이는 전력시장운영규칙 제2.4.2조의 증분가격(IP) 산정식과 대응된다:
\begin{equation}\label{eq:ip_rule}
    \text{IP}_{i,t} = \frac{(2 \times \text{QPC}_i \times \text{DAOS}_{i,t} + \text{LPC}_i)}{\text{TLF}} \times \frac{1}{1000}
\end{equation}
여기서 $\text{QPC}_i = a_g$, $\text{LPC}_i = b_g$, $\text{DAOS}_{i,t}$는 실제 운전점,
$\text{TLF}$는 송전손실계수(본 모형에서 1.0 가정)이다.

\textbf{중점 MC (Basic ED용)}: 발전기의 운전 범위 중점에서 산출하는 초기 추정치:
\begin{equation}\label{eq:mc_mid}
    \text{MC}_g^{\text{mid}} = \left(2a_g \cdot \frac{P_g^{\min} + P_g^{\max}}{2} + b_g\right) \times 1000 \qquad [\text{원/MWh}]
\end{equation}

\textbf{실제 운전점 MC}: ED 풀이 후 실제 발전량 $p_{g,t}$에서 재계산:
\begin{equation}\label{eq:mc_actual}
    \text{MC}_g^{\text{actual}}(t) = \begin{cases}
        (2a_g \cdot p_{g,t} + b_g) \times 1000 & \text{if } p_{g,t} > 10^{-3} \\
        (2a_g \cdot P_g^{\min} + b_g) \times 1000 & \text{if } p_{g,t} \leq 10^{-3}
    \end{cases}
\end{equation}

\subsection{SMP/수요 시계열}

시계열 데이터는 \texttt{smp\_demand.csv}에서 로딩되며, 각 레코드는 다음을 포함한다:
\begin{itemize}[noitemsep]
    \item \texttt{date}: 날짜
    \item \texttt{hour}: 시간 (1--24)
    \item \texttt{smp\_mainland}: 육지계통 SMP [원/MWh]
    \item \texttt{demand\_mainland}: 육지계통 수요 [MW]
\end{itemize}

\subsection{재생에너지 데이터}

재생에너지 발전량은 \texttt{renewables\_generation\_mwh.csv}에서 로딩된다.
2024년 데이터에서 발전량을 physical availability(potential)로 재구성하되,
peak 이용률(CF)이 1을 초과하지 않으므로 발전량 $\equiv$ potential로 가정한다:
\begin{equation}
    \text{potential}_{t}^{\text{PV}} = \text{solar\_gen}_t, \qquad
    \text{CF}_t^{\text{PV}} = \frac{\text{potential}_t^{\text{PV}}}{\text{Installed}_{\text{PV}}}
\end{equation}

\subsection{원전 Must-Off 스케줄}

원전 정비 일정은 \texttt{nuclear\_must\_off.csv}에서 로딩된다.
호기명(예: ``고리4호기'')을 발전기 ID(예: ``Nuclear\_003'')로 매핑하는 테이블을 사용하여,
해당 날짜에 정비 중인 원전의 출력을 $P_g^{\min} = P_g^{\max} = 0$으로 설정한다.

\subsection{연료비 및 호기 사양}

Price Adder의 물리적 상한을 검증하기 위해 \texttt{genthermal.csv}에서
고온기동비(startup cost) $S_g$ [천원], 최소가동시간 $\text{UT}_g$ [시간],
호기별 최대출력 $P_g^{\text{unit}}$ [MW]을 로딩한다.
물리적 상한은 다음과 같이 산출된다:
\begin{equation}\label{eq:adder_bound}
    \bar{A}_g = \frac{S_g \times 1000}{\text{UT}_g \times P_g^{\text{unit}}} \qquad [\text{원/MWh}]
\end{equation}


% ============================================================
% 3. Basic Economic Dispatch
% ============================================================
\section{Basic Economic Dispatch (Basic ED)}\label{sec:basic_ed}

\subsection{개요}

Basic ED는 가장 단순한 형태의 경제급전 모형으로,
재생에너지를 음의 부하로 처리하고 단순 선형비용으로 급전하는 벤치마크 모형이다.
최소출력, 램프 제약, must-run 조건이 없으므로 해석적으로 투명한 SMP를 도출한다.

\subsection{순수요 계산}

재생에너지는 음의 부하(negative load)로 처리하여 순수요를 산출한다:
\begin{equation}\label{eq:net_demand}
    D_t^{\text{net}} = D_t - \text{RE}_t \qquad \forall t \in \mathcal{T}
\end{equation}
순수요가 음수인 경우(재생에너지 $>$ 수요) 0으로 보정한다:
\begin{equation}
    D_t^{\text{net}} \leftarrow \max(0, D_t^{\text{net}})
\end{equation}

\subsection{수학적 정식화}

\subsubsection{목적함수 (B1)}

총 발전비용을 최소화한다:
\begin{equation}\label{eq:basic_obj}
    \boxed{\min_{p_{g,t}} \quad \sum_{t \in \mathcal{T}} \sum_{g \in \mathcal{G}} c_g \cdot p_{g,t}}
\end{equation}
여기서 $c_g$는 발전기 $g$의 한계비용 [\text{원/MWh}]으로, 시간에 따라 변하지 않는 상수이다.
구체적으로 $c_g = \text{MC}_g^{\text{mid}}$ (식 \eqref{eq:mc_mid})을 사용한다.

\subsubsection{수급균형 제약 (B2)}

각 시간대에서 총 발전량이 순수요와 일치해야 한다:
\begin{equation}\label{eq:basic_balance}
    \sum_{g \in \mathcal{G}} p_{g,t} = D_t^{\text{net}} \qquad \forall t \in \mathcal{T}
\end{equation}

\subsubsection{출력상한 제약 (B3)}

각 발전기의 출력은 비음이며 최대출력을 초과할 수 없다:
\begin{equation}\label{eq:basic_bounds}
    0 \leq p_{g,t} \leq P_g^{\max} \qquad \forall g \in \mathcal{G}, \; \forall t \in \mathcal{T}
\end{equation}

\subsubsection{SMP 도출 (B4)}

SMP는 수급균형 제약 \eqref{eq:basic_balance}의 쌍대변수(dual variable, shadow price)로 정의된다:
\begin{equation}\label{eq:basic_smp}
    \boxed{\lambda_t = \text{dual}\!\left(\sum_{g \in \mathcal{G}} p_{g,t} = D_t^{\text{net}}\right) \qquad \forall t \in \mathcal{T}}
\end{equation}
최소화 문제의 등호 제약에서 쌍대변수의 양수 값은
수요 1\,MW 증가 시 총비용 증가분(원/MWh)을 의미한다.

\subsection{한계연료원 식별}

각 시간대의 한계연료원(marginal fuel)은 부분 투입(partial dispatch) 상태인 발전기 중
비용이 가장 높은 발전기의 연료원으로 식별된다:
\begin{equation}
    g_t^* = \argmax_{g \in \mathcal{G}_t^{\text{marg}}} c_g
    , \qquad
    \mathcal{G}_t^{\text{marg}} = \left\{g : 10^{-3} < p_{g,t} < P_g^{\max} - 10^{-3}\right\}
\end{equation}
$\mathcal{G}_t^{\text{marg}} = \emptyset$이면 투입된 발전기 중 최고비용 발전기를 한계발전기로 간주한다.


% ============================================================
% 4. Pre-revision Economic Dispatch
% ============================================================
\section{Pre-revision Economic Dispatch (Pre-ED)}\label{sec:pre_ed}

\subsection{개요}

Pre-revision ED는 현행 한국 전력시장의 운영규칙을 반영한 모형으로,
Basic ED에 다음 요소를 추가한다:
\begin{itemize}[noitemsep]
    \item \textbf{유효 한계비용}: gencost 기반 MC + Price Adder
    \item \textbf{최소출력 제약}: must-run 발전기(원전)의 $P_g^{\min}$ 보장
    \item \textbf{램프 제약}: 시간간 출력 변동 제한
    \item \textbf{구간별 선형 비용 근사}: 2차 비용함수의 $S$개 구간 분할
    \item \textbf{출력제한(curtailment)}: 현행 must-run pmin 보전을 위한 RE 사전 cap
\end{itemize}

\subsection{유효 한계비용}

발전기 $g$의 시간 $t$에서의 유효 한계비용은 다음과 같이 구성된다:
\begin{equation}\label{eq:effective_mc}
    \tilde{c}_{g,t} = \text{MC}_g^{\text{eff}}(t) + A_{g,t}
\end{equation}
여기서:
\begin{itemize}[noitemsep]
    \item $\text{MC}_g^{\text{eff}}(t)$: gencost 기반 유효 한계비용 [원/MWh]
    \item $A_{g,t}$: Price Adder [원/MWh] --- calibration에서 추정(\Cref{sec:calibrate} 참조)
\end{itemize}

\subsection{출력제한 (Curtailment) 로직}

\subsubsection{현행 시장규칙 (\texttt{curtailment\_free = false})}

must-run 발전기의 최소출력 합 $P_{\text{must}}^{\min}$을 보전하기 위해
재생에너지를 사전에 cap한다:
\begin{align}
    P_{\text{must}}^{\min} &= \sum_{g \in \mathcal{G}_{\text{must}}} P_g^{\min} \\
    \text{RE}_t^{\text{eff}} &= \min\!\big(\text{RE}_t, \; D_t - P_{\text{must}}^{\min}\big) \label{eq:re_cap}
\end{align}
출력제한량: $\text{curt}_t = \text{RE}_t - \text{RE}_t^{\text{eff}}$

\subsubsection{Calibration 모드 (\texttt{curtailment\_free = true})}

Calibration 단계에서는 Price Adder가 출력제한의 dual pollution을
흡수하지 않도록 RE를 must-take 상수로 주입한다:
\begin{equation}
    \text{RE}_t^{\text{eff}} = \min(\text{RE}_t, D_t)
\end{equation}
이 설정은 Adder의 calibration 순도(purity)를 확보하기 위함이다.

\subsection{수학적 정식화 --- 단일 변수 방식}

\subsubsection{순수요}
\begin{equation}
    D_t^{\text{net}} = D_t - \text{RE}_t^{\text{eff}} \qquad \forall t \in \mathcal{T}
\end{equation}

\subsubsection{목적함수 (P2)}
\begin{equation}\label{eq:pre_obj}
    \boxed{\min_{p_{g,t}} \quad \sum_{t \in \mathcal{T}} \sum_{g \in \mathcal{G}} \tilde{c}_{g,t} \cdot p_{g,t}}
\end{equation}

\subsubsection{수급균형 제약 (P3)}
\begin{equation}\label{eq:pre_balance}
    \sum_{g \in \mathcal{G}} p_{g,t} = D_t^{\text{net}} \qquad \forall t \in \mathcal{T}
\end{equation}

\subsubsection{출력 제약 (P4)}

must-run 발전기는 최소출력을 보장하고, 비must-run 발전기는 0 이상의 출력을 갖는다:
\begin{equation}\label{eq:pre_bounds}
    \begin{cases}
        P_g^{\min} \leq p_{g,t} \leq P_g^{\max} & \text{if } g \in \mathcal{G}_{\text{must}} \\
        0 \leq p_{g,t} \leq P_g^{\max}           & \text{if } g \notin \mathcal{G}_{\text{must}}
    \end{cases}
    \qquad \forall g, t
\end{equation}

\subsubsection{램프 제약 (P5)}

연속 시간대 간의 출력 변동을 제한한다:
\begin{equation}\label{eq:ramp}
    -\text{RD}_g \leq p_{g,t} - p_{g,t-1} \leq \text{RU}_g \qquad \forall g \in \mathcal{G}, \; t \geq 2
\end{equation}
$\text{RU}_g = \infty$ 또는 $\text{RD}_g = \infty$인 경우 해당 방향의 제약을 생략한다.

\subsubsection{SMP 도출}
\begin{equation}
    \lambda_t = \text{dual}\!\left(\sum_{g \in \mathcal{G}} p_{g,t} = D_t^{\text{net}}\right) \qquad \forall t \in \mathcal{T}
\end{equation}

\subsection{수학적 정식화 --- 구간별 선형(Piecewise Linear) 방식}\label{sec:piecewise}

2차 비용함수를 $S$개 구간으로 선형 근사하여 LP 정식화를 유지한다.

\subsubsection{구간 분할}

발전기 $g$의 출력 범위 $[P_g^{\min}, P_g^{\max}]$를 $S$개의 등간격 구간으로 분할한다:
\begin{equation}\label{eq:pw_delta}
    \bar{\delta}_g = \frac{P_g^{\max} - P_g^{\min}}{S}
\end{equation}
각 구간 $s \in \{1, \ldots, S\}$의 중점:
\begin{equation}\label{eq:pw_mid}
    p_{\text{mid},s} = P_g^{\min} + \left(s - \frac{1}{2}\right) \cdot \bar{\delta}_g
\end{equation}
각 구간의 한계비용:
\begin{equation}\label{eq:pw_mc}
    \boxed{\text{mc}_{g,s} = (2a_g \cdot p_{\text{mid},s} + b_g) \times 1000 \qquad [\text{원/MWh}]}
\end{equation}

must-run이 아닌 발전기의 경우 $P_g^{\min}$을 0으로 설정하여
전체 범위 $[0, P_g^{\max}]$를 분할한다.

$|a_g| < 10^{-12}$인 선형 비용함수의 경우 단일 구간($S=1$)으로 처리:
$\text{mc}_{g,1} = b_g \times 1000$.

\subsubsection{증분 변수}

구간 $s$의 증분 변수 $\delta_{g,s,t} \geq 0$을 도입한다:
\begin{equation}\label{eq:pw_var}
    0 \leq \delta_{g,s,t} \leq \bar{\delta}_g \qquad \forall g, s, t
\end{equation}
총 출력은 기저점과 증분 변수의 합으로 표현된다:
\begin{equation}\label{eq:pw_total}
    p_{g,t}^{\text{total}} = P_g^{\min} + \sum_{s=1}^{S_g} \delta_{g,s,t}
\end{equation}
여기서 $S_g$는 발전기 $g$의 실제 세그먼트 수이다.

\subsubsection{Piecewise Linear 목적함수}

구간별 선형 근사를 사용하면 목적함수는 다음과 같이 변환된다:
\begin{equation}\label{eq:pw_obj}
    \boxed{\min \quad \underbrace{\sum_{t \in \mathcal{T}} \sum_{g \in \mathcal{G}} \sum_{s=1}^{S_g} \text{mc}_{g,s} \cdot \delta_{g,s,t}}_{\text{구간별 선형 비용}}
    + \underbrace{\sum_{t \in \mathcal{T}} \sum_{g \in \mathcal{G}} A_{g,t} \cdot p_{g,t}^{\text{total}}}_{\text{Price Adder 비용}}}
\end{equation}

\subsubsection{제약조건 (Piecewise 버전)}

수급균형, must-run 최소출력, 출력상한, 램프 제약은 모두
$p_{g,t}^{\text{total}}$ (식 \eqref{eq:pw_total})로 표현된다:
\begin{align}
    &\sum_{g \in \mathcal{G}} p_{g,t}^{\text{total}} = D_t^{\text{net}} && \forall t \label{eq:pw_balance} \\
    &p_{g,t}^{\text{total}} \geq P_g^{\min} && \forall g \in \mathcal{G}_{\text{must}}, t \label{eq:pw_must} \\
    &p_{g,t}^{\text{total}} \leq P_g^{\max} && \forall g, t \\
    &-\text{RD}_g \leq p_{g,t}^{\text{total}} - p_{g,t-1}^{\text{total}} \leq \text{RU}_g && \forall g, t \geq 2 \label{eq:pw_ramp}
\end{align}


% ============================================================
% 5. Post-revision Economic Dispatch
% ============================================================
\section{Post-revision Economic Dispatch (Post-ED)}\label{sec:post_ed}

\subsection{개요}

Post-revision ED는 재생에너지 입찰제 도입 후의 시장을 모형화한다.
Pre-ED 대비 다음 요소가 추가된다:
\begin{itemize}[noitemsep]
    \item \textbf{재생에너지 입찰 블록}: 6개 블록 (PV\_low, PV\_mid, PV\_high, W\_low, W\_mid, W\_high)
    \item \textbf{입찰 하한가(BidFloor)}: $-\beta \times \text{REC} \times 1000$
    \item \textbf{비입찰 RE 변수 분리}: dual pollution 방지를 위한 \texttt{re\_net} 변수
    \item \textbf{RE Pmin 제약}: 제주 규정 적용 --- 낙찰 시 최소 공급 의무
    \item \textbf{$\varepsilon_{\text{nonbid}}$ 페널티}: degeneracy 해소를 위한 비입찰 RE 비용항
\end{itemize}

\subsection{재생에너지 블록 구조}\label{sec:re_blocks}

\subsubsection{입찰참여분/비입찰분 분리}

전체 재생에너지 가용량을 입찰참여율 $\rho$에 따라 분리한다:
\begin{align}
    \text{PV}_t^{\text{bid}} &= \rho_{\text{PV}} \cdot \text{avail}_t^{\text{PV}} \label{eq:pv_bid} \\
    \text{W}_t^{\text{bid}} &= \rho_{\text{W}} \cdot \text{avail}_t^{\text{W}} \\
    \text{RE}_t^{\text{nonbid}} &= (1 - \rho_{\text{PV}}) \cdot \text{avail}_t^{\text{PV}}
        + (1 - \rho_{\text{W}}) \cdot \text{avail}_t^{\text{W}} \label{eq:re_nonbid}
\end{align}

\subsubsection{Low/Mid/High 블록 분할}

입찰참여분을 가중치 $\bm{w} = (w_{\text{low}}, w_{\text{mid}}, w_{\text{high}})$로 3분할한다
(태양광, 풍력 각각):
\begin{align}
    \bar{R}_{k,t}^{\text{PV\_low}}  &= w_{\text{low}}^{\text{PV}} \cdot \text{PV}_t^{\text{bid}} \\
    \bar{R}_{k,t}^{\text{PV\_mid}}  &= w_{\text{mid}}^{\text{PV}} \cdot \text{PV}_t^{\text{bid}} \\
    \bar{R}_{k,t}^{\text{PV\_high}} &= w_{\text{high}}^{\text{PV}} \cdot \text{PV}_t^{\text{bid}}
\end{align}
기본값: $\bm{w}^{\text{PV}} = \bm{w}^{\text{W}} = (0.4, 0.3, 0.3)$.
총 6개 블록: $\mathcal{K} = \{\text{PV\_low}, \text{PV\_mid}, \text{PV\_high}, \text{W\_low}, \text{W\_mid}, \text{W\_high}\}$.

\subsubsection{블록별 설비용량}

각 블록이 점유하는 설비용량은 다음과 같이 산출된다:
\begin{equation}\label{eq:block_inst}
    \text{inst}_k = w_k \cdot \rho \cdot \text{Installed}_{\text{tech}}
\end{equation}
여기서 $\text{tech} \in \{\text{PV}, \text{W}\}$.

\subsubsection{입찰 하한가 (BidFloor)}

제주형 하한가 규정을 적용한다:
\begin{equation}\label{eq:bid_floor}
    \boxed{\text{BidFloor} = -\beta \times \text{REC} \times 1000 \qquad [\text{원/MWh}]}
\end{equation}
여기서:
\begin{itemize}[noitemsep]
    \item $\beta$: 하한가 계수 (기본값 2.0, 민감도 분석 시 1.5/2.0/2.5)
    \item $\text{REC}$: REC 평균가격 [원/kWh] (기본값 80 원/kWh)
    \item $\times 1000$: kWh $\to$ MWh 단위 변환
\end{itemize}
예: $\beta = 2.0$, $\text{REC} = 80$ 원/kWh $\Rightarrow$ BidFloor $= -160{,}000$ 원/MWh.

\subsubsection{시나리오별 입찰가격 부여}

4가지 입찰 시나리오에 따른 블록별 입찰가격:

\begin{table}[H]
\centering
\caption{시나리오별 입찰가격 ($f$: BidFloor)}\label{tab:bid_scenarios}
\small
\begin{tabular}{lcccc}
\toprule
\textbf{블록} & \textbf{zero} & \textbf{floor} & \textbf{mixed} & \textbf{conservative} \\
\midrule
Low  & 0 & $f$ & $f$    & $0.5f$ \\
Mid  & 0 & $f$ & $0.5f$ & $0.25f$ \\
High & 0 & $f$ & 0      & 0 \\
\bottomrule
\end{tabular}
\end{table}

\subsection{수학적 정식화}

\subsubsection{결정변수}

\begin{itemize}[noitemsep]
    \item $p_{g,t}$: 열발전기 $g$의 시간 $t$ 발전량 [MW] (또는 구간별 $\delta_{g,s,t}$)
    \item $r_{k,t}$: 재생에너지 입찰 블록 $k$의 시간 $t$ 낙찰량 [MW]
    \item $\text{re\_net}_t$: 비입찰 재생에너지의 실제 투입량 [MW]
\end{itemize}

\subsubsection{목적함수 (R1)}

\textbf{Piecewise Linear 버전}:
\begin{equation}\label{eq:post_obj}
    \boxed{\min \quad
    \underbrace{\sum_{t,g,s} \text{mc}_{g,s} \cdot \delta_{g,s,t}}_{\text{구간별 선형 비용}}
    + \underbrace{\sum_{t,g} A_{g,t} \cdot p_{g,t}^{\text{total}}}_{\text{Price Adder}}
    + \underbrace{\sum_{t,k} b_{k,t} \cdot r_{k,t}}_{\text{RE 입찰 비용}}
    + \underbrace{\varepsilon_{\text{nonbid}} \cdot \sum_t \text{re\_net}_t}_{\text{비입찰 RE 비용}}}
\end{equation}
여기서:
\begin{itemize}[noitemsep]
    \item $\text{mc}_{g,s}$: 구간 $s$의 한계비용 (식 \eqref{eq:pw_mc})
    \item $A_{g,t}$: Price Adder
    \item $b_{k,t}$: 블록 $k$의 입찰가격
    \item $\varepsilon_{\text{nonbid}} = 100$ 원/MWh: degeneracy 해소용 소규모 비용
\end{itemize}

\textbf{단일 변수 버전} (Piecewise 미사용 시):
\begin{equation}
    \min \quad \sum_{t,g} \tilde{c}_{g,t} \cdot p_{g,t}
    + \sum_{t,k} b_{k,t} \cdot r_{k,t}
    + \varepsilon_{\text{nonbid}} \sum_t \text{re\_net}_t
\end{equation}

\paragraph{$\varepsilon_{\text{nonbid}}$ 항의 역할.}
비입찰 RE($\text{re\_net}_t$)에 소규모 양의 비용을 부여함으로써,
LP 솔버가 비입찰 RE와 입찰 RE 간의 교환에서 degeneracy를 겪지 않도록 한다.
이 항이 없으면 dual(SMP)이 불안정하게 결정될 수 있다.

\subsubsection{수급균형 제약 (R2)}

총 수요를 열발전, 재생 입찰블록, 비입찰 RE가 충족한다:
\begin{equation}\label{eq:post_balance}
    \boxed{\sum_{g \in \mathcal{G}} p_{g,t}^{\text{total}} + \sum_{k \in \mathcal{K}} r_{k,t} + \text{re\_net}_t = D_t \qquad \forall t \in \mathcal{T}}
\end{equation}

\textbf{주의}: Pre-ED와 달리 순수요가 아닌 \emph{총수요} $D_t$를 사용한다.
비입찰 RE는 $\text{re\_net}_t$ 변수로 내생적으로 결정된다.

\subsubsection{비입찰 RE 범위}

dual pollution 방지를 위해 $\text{re\_net}_t$를 제한된 범위의 변수로 정의한다:
\begin{equation}\label{eq:re_net_bound}
    0 \leq \text{re\_net}_t \leq \text{RE}_t^{\text{nonbid}} \qquad \forall t \in \mathcal{T}
\end{equation}
출력제한량은 사후 계산: $\text{curt}_t = \text{RE}_t^{\text{nonbid}} - \text{re\_net}_t$.

\subsubsection{RE 입찰 블록 제약 (R3)}

각 입찰 블록의 낙찰량은 공급가능량 이하이며, RE Pmin 제약을 만족해야 한다:
\begin{equation}\label{eq:re_bounds}
    P_{k,t}^{\min,\text{RE}} \leq r_{k,t} \leq \bar{R}_{k,t} \qquad \forall k \in \mathcal{K}, \; \forall t \in \mathcal{T}
\end{equation}

\subsubsection{RE Pmin 제약 (제주 규정)}

제주 전력시장 규정을 참고하여, 입찰 블록이 낙찰될 경우 최소 공급 의무를 부과한다:
\begin{equation}\label{eq:re_pmin}
    \boxed{P_{k,t}^{\min,\text{RE}} = \begin{cases}
        \min(\alpha \cdot \text{inst}_k, \; \bar{R}_{k,t}) & \text{if } \text{inst}_k > 0 \text{ and } \bar{R}_{k,t} > 10^{-3} \\
        \alpha \cdot \bar{R}_{k,t} & \text{if } \text{inst}_k = 0 \text{ and } \bar{R}_{k,t} > 10^{-3} \\
        0 & \text{otherwise}
    \end{cases}}
\end{equation}
여기서 $\alpha$ = \texttt{re\_pmin\_frac} (기본값 0.1).

\subsubsection{열발전 제약}

Pre-ED의 모든 제약 (P4), (P5)를 그대로 유지한다:
\begin{align}
    &\text{(P4)} \quad
    \begin{cases}
        P_g^{\min} \leq p_{g,t}^{\text{total}} \leq P_g^{\max} & g \in \mathcal{G}_{\text{must}} \\
        0 \leq p_{g,t}^{\text{total}} \leq P_g^{\max}           & g \notin \mathcal{G}_{\text{must}}
    \end{cases} \\
    &\text{(P5)} \quad
    -\text{RD}_g \leq p_{g,t}^{\text{total}} - p_{g,t-1}^{\text{total}} \leq \text{RU}_g \qquad t \geq 2
\end{align}

\subsubsection{SMP 도출}

$\text{re\_net}$ 변수 분리에 의해 출력제한 페널티가 dual에 포함되지 않으므로,
LP dual을 직접 SMP로 사용한다:
\begin{equation}\label{eq:post_smp}
    \boxed{\lambda_t = \text{dual}\!\left(\sum_g p_{g,t}^{\text{total}} + \sum_k r_{k,t} + \text{re\_net}_t = D_t\right)}
\end{equation}

\subsection{Case A: Zero-Bidding Baseline ($\rho = 0$)}

Case A (Case\_A\_zero)에서는 입찰참여율 $\rho_{\text{PV}} = \rho_{\text{W}} = 0$으로 설정하여
재생에너지 입찰 블록 변수를 생성하지 않는다 (\texttt{bidding\_active = false}).
이 경우 Post-ED는 Pre-ED와 동치가 되어야 하며, 이를 검증한다:
\begin{equation}
    \max_t |\text{SMP}_t^{\text{post,A}} - \text{SMP}_t^{\text{pre}}| < 1 \; \text{원/MWh}
\end{equation}

\subsection{$\Delta$SMP 분석}

Pre vs Post의 SMP 변화를 분석한다:
\begin{equation}\label{eq:delta_smp}
    \Delta\text{SMP}_t = \text{SMP}_t^{\text{post}} - \text{SMP}_t^{\text{pre}}
\end{equation}
\begin{itemize}[noitemsep]
    \item $\Delta\text{SMP}_t > 0$: SMP 상승 (저녁 램프 등)
    \item $\Delta\text{SMP}_t < 0$: SMP 하락 (낮 시간 재생 주입 등)
\end{itemize}


% ============================================================
% 6. Price Adder 보정
% ============================================================
\section{Price Adder 보정 (Calibration)}\label{sec:calibrate}

\subsection{동기}

경제급전 모형은 기동비(startup cost), 무부하비(no-load cost), 정지비(shutdown cost) 등
기동정지(UC, Unit Commitment) 관련 비용을 직접 모형화하지 않으므로,
실제 SMP와 체계적 괴리가 발생한다.
Price Adder $A_{g,t}$는 이러한 미모형 요소를 부분 흡수하는 보정항으로,
반복 보정법(iterative correction)으로 추정한다.

\subsection{검증지표}

\subsubsection{MAE (Mean Absolute Error)}
\begin{equation}\label{eq:mae}
    \text{MAE} = \frac{1}{T} \sum_{t=1}^{T} |\text{SMP}_t^{\text{model}} - \text{SMP}_t^{\text{actual}}|
\end{equation}

\subsubsection{RMSE (Root Mean Square Error)}
\begin{equation}\label{eq:rmse}
    \text{RMSE} = \sqrt{\frac{1}{T} \sum_{t=1}^{T} (\text{SMP}_t^{\text{model}} - \text{SMP}_t^{\text{actual}})^2}
\end{equation}

\subsubsection{지속곡선(Duration Curve) 오차}
\begin{equation}
    \text{DC\_Error} = \frac{1}{T} \sum_{t=1}^{T} |\text{DC}_t^{\text{model}} - \text{DC}_t^{\text{actual}}|
\end{equation}
여기서 $\text{DC}_t$는 SMP를 내림차순 정렬한 지속곡선이다.

\subsection{단일일 Price Adder 추정}\label{sec:single_day_adder}

\subsubsection{알고리즘 개요}

$A_{g,t}$를 0으로 초기화한 후, 각 반복에서:
\begin{enumerate}[noitemsep]
    \item Pre-ED를 풀어 $\text{SMP}_t^{\text{model}}$을 산출
    \item 오차 $e_t = \text{SMP}_t^{\text{actual}} - \text{SMP}_t^{\text{model}}$ 계산
    \item 활성 한계발전기(marginal generator) 집합 식별
    \item Adder 갱신
    \item L2 shrinkage 적용
    \item 물리적 bounds clamping
\end{enumerate}

\subsubsection{활성 한계발전기 집합}

시간 $t$에서 부분 투입 상태인 발전기 집합:
\begin{equation}\label{eq:marg_set}
    \mathcal{M}_t = \left\{g \in \mathcal{G} : P_g^{\min} + 10^{-3} < p_{g,t} < P_g^{\max} - 10^{-3}\right\}
\end{equation}
여기서 must-run이 아닌 발전기의 $P_g^{\min}$은 0으로 간주한다.

\subsubsection{Adder 갱신 규칙 (1/$n_{\text{marg}}$ 정규화)}

\begin{equation}\label{eq:adder_update}
    \boxed{A_{g,t} \leftarrow A_{g,t} + \eta \cdot \frac{e_t}{n_{\text{marg},t}} \qquad \forall g \in \mathcal{M}_t}
\end{equation}
여기서:
\begin{itemize}[noitemsep]
    \item $\eta$: 학습률 (기본값 0.3)
    \item $e_t = \text{SMP}_t^{\text{actual}} - \text{SMP}_t^{\text{model}}$: 시간 $t$의 SMP 오차
    \item $n_{\text{marg},t} = |\mathcal{M}_t|$: 활성 한계발전기 수
\end{itemize}

$n_{\text{marg},t} = 0$이면 갱신을 생략한다.
$1/n_{\text{marg}}$ 정규화는 오차를 한계발전기들에 균등 분배하여
단일 발전기에 과도한 보정이 집중되는 것을 방지한다.

\subsubsection{L2 Shrinkage (Tikhonov 정규화)}

각 반복 후 Adder 전체에 shrinkage를 적용한다:
\begin{equation}\label{eq:l2_shrinkage}
    \boxed{A_{g,t} \leftarrow (1 - \lambda\eta) \cdot A_{g,t} \qquad \forall g, t}
\end{equation}
여기서 $\lambda\eta$ = \texttt{l2\_shrinkage} (기본값 0.05).

이는 Tikhonov 정규화의 경사하강법 등가로,
Adder가 0으로 수축하는 효과를 제공하여 과적합을 방지한다.

\subsubsection{물리적 Bounds Clamping}

Adder가 기동비 기반 물리적 상한(식 \eqref{eq:adder_bound})을 초과하지 않도록 제한한다:
\begin{equation}\label{eq:clamp}
    A_{g,t} \leftarrow \text{clamp}(A_{g,t}, -\bar{A}_g, \bar{A}_g)
\end{equation}

\subsubsection{수렴 판정}

MAE가 목표값(기본값 5,000 원/MWh) 이하이면 조기 종료한다.

\begin{algorithm}[H]
\caption{단일일 Price Adder 추정}\label{alg:single_adder}
\begin{algorithmic}[1]
\REQUIRE $\mathcal{G}$, $\mathcal{T}$, $\text{SMP}^{\text{actual}}$, $\eta$, $\lambda\eta$, $\bar{A}$, $N_{\max}$
\STATE $A_{g,t} \leftarrow 0 \quad \forall g, t$
\FOR{$\text{iter} = 1$ to $N_{\max}$}
    \STATE Pre-ED를 풀어 $\text{SMP}^{\text{model}}$ 산출 (\texttt{curtailment\_free = true})
    \STATE $\text{MAE} \leftarrow \frac{1}{T}\sum_t |e_t|$ \quad 여기서 $e_t = \text{SMP}_t^{\text{actual}} - \text{SMP}_t^{\text{model}}$
    \IF{$\text{MAE} < \text{target\_mae}$}
        \STATE \textbf{break}
    \ENDIF
    \FOR{$t = 1$ to $T$}
        \STATE $\mathcal{M}_t \leftarrow \{g : P_g^{\min} + \epsilon < p_{g,t} < P_g^{\max} - \epsilon\}$
        \STATE $n_{\text{marg}} \leftarrow |\mathcal{M}_t|$
        \IF{$n_{\text{marg}} > 0$}
            \FOR{$g \in \mathcal{M}_t$}
                \STATE $A_{g,t} \leftarrow A_{g,t} + \eta \cdot e_t / n_{\text{marg}}$
            \ENDFOR
        \ENDIF
    \ENDFOR
    \STATE $A \leftarrow (1 - \lambda\eta) \cdot A$ \hfill \textit{// L2 shrinkage}
    \STATE $A_{g,t} \leftarrow \text{clamp}(A_{g,t}, -\bar{A}_g, \bar{A}_g) \quad \forall g, t$ \hfill \textit{// bounds}
\ENDFOR
\RETURN $A$
\end{algorithmic}
\end{algorithm}

\subsection{다일(Multi-day) 3D Adder 추정}\label{sec:multi_day_adder}

단일일 Adder는 특정 일자에 과적합될 수 있으므로,
다수의 훈련일에 걸쳐 계절별 3차원 Adder를 추정한다.

\subsubsection{3D Adder 구조}

\begin{equation}\label{eq:adder_3d}
    \boxed{A[g, h, s] \in \mathbb{R}^{G \times 24 \times 4}}
\end{equation}
여기서:
\begin{itemize}[noitemsep]
    \item $g \in \{1, \ldots, G\}$: 발전기 인덱스
    \item $h \in \{1, \ldots, 24\}$: 시간 인덱스
    \item $s \in \{1, 2, 3, 4\}$: 계절 인덱스 (봄/여름/가을/겨울)
\end{itemize}

계절 분류:
\begin{equation}
    s = \begin{cases}
        1 \; (\text{spring}) & \text{월} \in \{3, 4, 5\} \\
        2 \; (\text{summer}) & \text{월} \in \{6, 7, 8\} \\
        3 \; (\text{fall})   & \text{월} \in \{9, 10, 11\} \\
        4 \; (\text{winter}) & \text{월} \in \{12, 1, 2\}
    \end{cases}
\end{equation}

\subsubsection{Train/Test/Buffer 분할}

시간적 누출(temporal leakage)을 방지하기 위해 다음과 같이 분할한다:
\begin{itemize}[noitemsep]
    \item \textbf{Test}: 계절당 3일, 총 12일 (대표일)
    \item \textbf{Buffer}: Test 일자 $\pm 3$일 (제외 구간)
    \item \textbf{Train}: 계절당 25일, 총 $\approx$100일 (계절균형 층화추출)
\end{itemize}

\subsubsection{에폭 단위 갱신}

각 에폭(epoch)에서 훈련일을 무작위 순서로 순회하며 누적 갱신을 수행한다:
\begin{equation}
    \Delta A[g, h, s] = \frac{1}{n_s} \sum_{d \in \mathcal{D}_s^{\text{train}}} \eta \cdot \frac{e_{d,h}}{n_{\text{marg},d,h}}
    \cdot \mathbb{1}[g \in \mathcal{M}_{d,h}]
\end{equation}
여기서 $n_s$는 계절 $s$에 속하는 훈련일 수이다.

에폭 종료 시:
\begin{align}
    A[g, h, s] &\leftarrow A[g, h, s] + \Delta A[g, h, s] && \text{(계절별 평균 갱신)} \\
    A &\leftarrow (1 - \lambda\eta) \cdot A && \text{(L2 shrinkage)} \\
    A[g, h, s] &\leftarrow \text{clamp}(A[g, h, s], -\bar{A}_g, \bar{A}_g) && \text{(물리적 bounds)}
\end{align}

\subsubsection{일자별 Adder 슬라이스}

특정 날짜 $d$에 대한 Adder는 해당 계절의 슬라이스를 사용한다:
\begin{equation}
    A_{g,t}(d) = A[g, t, s(d)]
\end{equation}
여기서 $s(d)$는 날짜 $d$의 계절 인덱스이다.

\subsection{Curtailment-free Calibration Purity}

Calibration 과정에서 \texttt{curtailment\_free = true}로 설정하여
RE를 must-take로 주입한다. 이는 출력제한의 dual이 Adder에 유입되어
Adder의 물리적 해석을 오염시키는 것(dual pollution)을 방지한다.

\subsection{교차검증}

Leave-one-out 교차검증으로 과적합 여부를 평가한다:
\begin{equation}
    R_{\text{overfit}} = \frac{\text{MAE}_{\text{test}}}{\text{MAE}_{\text{train}}}
\end{equation}
$R_{\text{overfit}} > 1.5$이면 과적합을 의심한다.


% ============================================================
% 7. 시나리오 분석
% ============================================================
\section{시나리오 분석}\label{sec:scenarios}

\subsection{4개 기본 시나리오}

\begin{table}[H]
\centering
\caption{4개 기본 시나리오 (Case A--D)}\label{tab:scenarios}
\begin{tabularx}{\textwidth}{lllXcc}
\toprule
\textbf{시나리오} & \textbf{코드명} & \textbf{입찰모드} & \textbf{설명} & $\rho_{\text{PV}}$ & $\rho_{\text{W}}$ \\
\midrule
Case A & Case\_A\_zero         & zero         & 입찰참여 없음 ($\rho=0$), baseline        & 0.0 & 0.0 \\
Case B & Case\_B\_floor        & floor        & 전 블록 하한가 입찰                        & 0.3 & 0.3 \\
Case C & Case\_C\_mixed        & mixed        & Low=하한가, Mid=50\%, High=0               & 0.3 & 0.3 \\
Case D & Case\_D\_conservative & conservative & Low=50\%, Mid=25\%, High=0                 & 0.3 & 0.3 \\
\bottomrule
\end{tabularx}
\end{table}

Case A는 $\rho = 0$이므로 입찰 블록 변수를 생성하지 않으며(\texttt{bidding\_active = false}),
Post-ED와 Pre-ED의 SMP가 동치임을 검증하는 baseline 역할을 한다.

\subsection{$\beta$ 민감도 분석}

입찰 하한가 계수 $\beta \in \{1.5, 2.0, 2.5\}$를 변동하여 SMP 영향을 분석한다.
$\beta$가 클수록 음의 입찰가격(BidFloor)이 더 낮아져 재생에너지의 가격 하방 압력이 강해진다:
\begin{equation}
    |\text{BidFloor}| = \beta \times 80 \times 1000 \qquad [\text{원/MWh}]
\end{equation}

\subsection{$\rho$ 민감도 분석}

입찰참여율 $\rho \in \{0.1, 0.2, 0.3, 0.5\}$를 변동하여 SMP 영향을 분석한다.
$\rho$가 높을수록 더 많은 재생에너지가 입찰에 참여하여 공급곡선에 진입한다.

\subsection{BidderType 이질성 (Heterogeneity)}\label{sec:bidder_types}

재생에너지 입찰사업자를 4가지 유형으로 분류하여 행태 이질성을 반영한다:

\begin{table}[H]
\centering
\caption{입찰사업자 유형}\label{tab:bidder_types}
\small
\begin{tabular}{lccccc}
\toprule
\textbf{유형} & $\phi_j$ (비율) & Beta$(\alpha, \beta)$ & $w_{\text{low}}$ & $w_{\text{mid}}$ & $w_{\text{high}}$ \\
\midrule
aggressive   & 0.30 & Beta(2, 5) & 0.6 & 0.3 & 0.1 \\
moderate     & 0.40 & Beta(3, 3) & 0.4 & 0.4 & 0.2 \\
conservative & 0.20 & Beta(5, 2) & 0.2 & 0.4 & 0.4 \\
PPA\_locked  & 0.10 & Beta(8, 1.5) & 0.1 & 0.3 & 0.6 \\
\bottomrule
\end{tabular}
\end{table}

유형 비율 $\sum_j \phi_j = 1$을 만족하며, 각 유형의 Beta 분포는
입찰가격 결정 시 확률적 행태를 모형화한다.

\subsubsection{정책 침투도 시나리오}

시장 성숙도에 따른 사업자 유형 구성 변화:

\begin{table}[H]
\centering
\caption{정책 침투도 시나리오별 유형 비율}\label{tab:penetration}
\small
\begin{tabular}{lcccc}
\toprule
\textbf{시나리오} & aggressive & moderate & conservative & PPA\_locked \\
\midrule
S1\_Early      & 0.50 & 0.30 & 0.15 & 0.05 \\
S2\_Mature     & 0.30 & 0.40 & 0.20 & 0.10 \\
S3\_Aggressive & 0.15 & 0.35 & 0.30 & 0.20 \\
\bottomrule
\end{tabular}
\end{table}

\subsection{Beta 혼합 + Common Shock 몬테카를로}\label{sec:monte_carlo}

\subsubsection{확률 모형}

각 시뮬레이션 표본 $i$에서:

\paragraph{Step 1: 유형별 확률 변수 추출.}
각 사업자 유형 $j$에 대해 Beta 분포에서 표본을 추출한다:
\begin{equation}\label{eq:beta_draw}
    u_j \sim \text{Beta}(\alpha_j, \beta_j) \qquad j = 1, \ldots, J
\end{equation}

\paragraph{Step 2: Common Shock 추출.}
시장 전체에 영향을 미치는 공통 충격:
\begin{equation}\label{eq:common_shock}
    \kappa \sim \max\!\big(0, \; \mathcal{N}(1, \sigma^2)\big)
\end{equation}
기본값: $\sigma = 0.10$. $\kappa$를 0 이상으로 절단한다.

\paragraph{Step 3: 블록별 입찰가 결정.}
Low/Mid/High 각 블록 $l$에 대해, 유형별 가중평균으로 통합 입찰 강도 $\bar{u}_l$을 산출한다:
\begin{equation}\label{eq:u_block}
    \bar{u}_l = \frac{\sum_{j=1}^{J} \phi_j \cdot w_{j,l} \cdot u_j}{\sum_{j=1}^{J} \phi_j \cdot w_{j,l}}
    \qquad l \in \{\text{low}, \text{mid}, \text{high}\}
\end{equation}
여기서 $w_{j,l}$은 유형 $j$의 블록 $l$ 가중치이다.

입찰가격은 다음과 같이 결정된다:
\begin{equation}\label{eq:bid_price_mc}
    \boxed{b_l = \text{clamp}\!\big(\kappa \cdot (\bar{u}_l - 1) \cdot |\text{BidFloor}|, \; \text{BidFloor}, \; 0\big)}
\end{equation}
여기서 $|\text{BidFloor}| = \beta \times \text{REC} \times 1000$.

\subsubsection{시뮬레이션 실행}

$N$ 표본(기본값 200)에 대해 Post-ED를 반복 실행하고 결과를 집계한다:
\begin{align}
    \overline{\text{SMP}}_t &= \frac{1}{N_{\text{success}}} \sum_{i=1}^{N_{\text{success}}} \text{SMP}_t^{(i)} \\
    \text{SMP}_t^{5\%} &= Q_{0.05}\!\left(\{\text{SMP}_t^{(i)}\}_{i=1}^{N_{\text{success}}}\right) \\
    \text{SMP}_t^{95\%} &= Q_{0.95}\!\left(\{\text{SMP}_t^{(i)}\}_{i=1}^{N_{\text{success}}}\right) \\
    \overline{\Delta\text{SMP}} &= \frac{1}{T} \sum_{t=1}^{T} \left(\overline{\text{SMP}}_t - \text{SMP}_t^{\text{pre}}\right)
\end{align}

\subsection{출력제한 분석}

각 시나리오의 출력제한을 분석한다:
\begin{align}
    \text{Curt}_{\text{total}} &= \sum_{t=1}^{T} \text{curt}_t \qquad [\text{MWh}] \\
    \text{Curt}_{\text{hours}} &= |\{t : \text{curt}_t > 10^{-3}\}| \\
    \text{Corr}(\text{curt}, \text{SMP}) &= \frac{\text{Cov}(\text{curt}_t, \text{SMP}_t)}
        {\text{Std}(\text{curt}_t) \cdot \text{Std}(\text{SMP}_t)}
\end{align}

Pre vs Post의 출력제한 감소율:
\begin{equation}
    \text{Reduction}\,\% = \left(1 - \frac{\text{Curt}_{\text{post}}}{\text{Curt}_{\text{pre}}}\right) \times 100
\end{equation}


% ============================================================
% 8. 구현 세부사항
% ============================================================
\section{구현 세부사항}\label{sec:implementation}

\subsection{소프트웨어 환경}

\begin{table}[H]
\centering
\caption{소프트웨어 스택}\label{tab:software}
\begin{tabular}{ll}
\toprule
\textbf{구성요소} & \textbf{설명} \\
\midrule
프로그래밍 언어 & Julia \\
수리 최적화     & JuMP (Julia for Mathematical Programming) \\
LP 솔버        & HiGHS (High Performance Optimization Software) \\
데이터 처리     & DataFrames.jl, CSV.jl \\
통계 분포       & Distributions.jl (Beta, Normal) \\
난수 생성       & Random (MersenneTwister, seed=2024) \\
\bottomrule
\end{tabular}
\end{table}

\subsection{소스코드 구조}

\begin{table}[H]
\centering
\caption{소스코드 파일 구조}\label{tab:source_files}
\small
\begin{tabularx}{\textwidth}{lX}
\toprule
\textbf{파일} & \textbf{역할} \\
\midrule
\texttt{types.jl}          & 핵심 자료형 정의 (ThermalGenerator, RenewableBidBlock, BidderType 등) \\
\texttt{load\_data.jl}     & CSV 데이터 로딩 (generators, gencost, SMP/수요, 원전 정비 등) \\
\texttt{preprocess.jl}     & 대표일 선정, MC 계산, Nuclear Must-Off, Piecewise Linear 구간 분할 \\
\texttt{build\_basic\_ed.jl} & Basic ED 모형 (순수 LP) \\
\texttt{build\_pre\_ed.jl}   & Pre-revision ED 모형 (유효MC, 램프, must-run, PW) \\
\texttt{build\_post\_ed.jl}  & Post-revision ED 모형 (6블록, re\_net, BidFloor, RE Pmin) \\
\texttt{calibrate.jl}       & Price Adder 추정 (단일일/다일, L2 shrinkage, 교차검증) \\
\texttt{scenarios.jl}       & 시나리오 실행, 민감도, Beta 몬테카를로 \\
\texttt{run\_all.jl}        & 전체 파이프라인 오케스트레이션 \\
\texttt{verify\_blocks.jl}  & RE 블록 생성 및 SMP 결정 검증 \\
\bottomrule
\end{tabularx}
\end{table}

\subsection{파이프라인 구조}

전체 분석 파이프라인은 6단계(PHASE)로 구성된다:

\begin{description}[style=nextline, noitemsep]
    \item[PHASE 0: 데이터 로딩]
    CSV 파일에서 8,784시간 통합 패널을 빌드한다.
    발전기 122기, gencost 계수, 원전 정비 일정, RE 설비용량을 로딩하고,
    Piecewise Linear 비용함수($S=4$)를 사전 계산한다.

    \item[PHASE 1: Train/Test/Buffer 분할]
    일별 프로파일(수요 피크, 평균 SMP, RE 비율, 저녁 램프)을 계산하여
    계절당 3개 대표일(test 12일), buffer($\pm$3일), train($\approx$100일)을 분할한다.

    \item[PHASE 2: Multi-day Calibration]
    3D Adder $A[G \times 24 \times 4]$를 10 에폭에 걸쳐 추정한다.
    학습률 $\eta = 0.2$, L2 shrinkage $\lambda\eta = 0.05$, target MAE $= 4{,}000$ 원/MWh.

    \item[PHASE 3: 대표일 평가]
    12개 test일에서 Pre-ED/Post-ED를 실행하고 4개 시나리오(Case A--D)를 평가한다.
    각 일자에서 원전 Must-Off를 적용하고, PW cost를 재계산한다.

    \item[PHASE 4: 정책 침투도 시나리오]
    S1\_Early, S2\_Mature, S3\_Aggressive에 대해 Beta 몬테카를로(100 표본)를 실행한다.

    \item[PHASE 5: 민감도 분석]
    $\beta \in \{1.5, 2.0, 2.5\}$, $\rho \in \{0.1, 0.2, 0.3, 0.5\}$에 대해 시나리오를 실행한다.
\end{description}

\subsection{JuMP 모형 구현 상세}

\subsubsection{모형 생성}

각 ED 모형은 다음 패턴으로 구현된다:
\begin{enumerate}[noitemsep]
    \item \texttt{Model(HiGHS.Optimizer)} --- 솔버 인스턴스 생성
    \item \texttt{set\_silent(model)} --- 솔버 출력 억제
    \item \texttt{@variable} --- 결정변수 선언 (bounds 포함)
    \item \texttt{@constraint} --- 수급균형, 램프, must-run 제약
    \item \texttt{@objective} --- 목적함수 설정 (Min)
    \item \texttt{optimize!(model)} --- 최적화 실행
    \item \texttt{dual(balance[t])} --- SMP 추출
\end{enumerate}

\subsubsection{Piecewise Linear 구현}

구간별 증분 변수 \texttt{delta[g, s, t]}를 3차원 배열로 선언하며,
사용하지 않는 세그먼트(발전기의 $S_g < S_{\max}$)는 \texttt{fix(delta, 0.0)}으로 고정한다.
총 출력 \texttt{p\_total[g, t]}은 \texttt{@expression}으로 정의하여
목적함수와 제약조건에서 재사용한다.

\subsubsection{수치적 고려사항}

\begin{itemize}[noitemsep]
    \item 부분 투입 판별 임계값: $\epsilon = 10^{-3}$ MW
    \item LP dual 부호 관례: JuMP의 Min 문제에서 등호 제약의 dual은 양수
    \item Piecewise Linear: 구간 한계비용이 단조증가하므로 LP 완화(relaxation)가 정확
    \item MathOptInterface 상태 확인: \texttt{MOI.OPTIMAL} 여부 검증
\end{itemize}

\subsection{데이터 단위 체계}

\begin{table}[H]
\centering
\caption{주요 변수의 단위}\label{tab:units}
\begin{tabular}{lll}
\toprule
\textbf{변수} & \textbf{단위} & \textbf{비고} \\
\midrule
발전출력 $p_{g,t}$     & MW     & \\
수요 $D_t$             & MW     & \\
SMP $\lambda_t$        & 원/MWh & \\
gencost $C(P)$         & 천원/h & $a$: 천원/MW$^2$h, $b$: 천원/MWh \\
한계비용 MC             & 원/MWh & gencost에서 $\times 1000$ 변환 \\
BidFloor               & 원/MWh & REC(원/kWh) $\times 1000$ \\
Price Adder $A_{g,t}$  & 원/MWh & \\
기동비 $S_g$           & 천원   & 물리적 상한 산출 시 $\times 1000$ \\
REC 가격               & 원/kWh & 기본값 80 원/kWh \\
\bottomrule
\end{tabular}
\end{table}


% ============================================================
% 9. 수학적 정식화 요약
% ============================================================
\section{수학적 정식화 종합 요약}\label{sec:summary}

본 절에서는 3단계 ED 모형의 수학적 정식화를 종합 비교한다.

\subsection{Basic ED}

\begin{align}
    &\min_{p_{g,t}} && \sum_{t \in \mathcal{T}} \sum_{g \in \mathcal{G}} c_g \cdot p_{g,t} \tag{B1} \\
    &\text{s.t.} && \sum_{g \in \mathcal{G}} p_{g,t} = D_t - \text{RE}_t && \forall t \tag{B2} \\
    & && 0 \leq p_{g,t} \leq P_g^{\max} && \forall g, t \tag{B3}
\end{align}

\subsection{Pre-revision ED (Piecewise Linear)}

\begin{align}
    &\min_{\delta_{g,s,t}} && \sum_{t,g,s} \text{mc}_{g,s} \cdot \delta_{g,s,t}
        + \sum_{t,g} A_{g,t} \cdot p_{g,t}^{\text{total}} \tag{P2'} \\
    &\text{s.t.} && \sum_{g \in \mathcal{G}} p_{g,t}^{\text{total}} = D_t - \text{RE}_t^{\text{eff}} && \forall t \tag{P3} \\
    & && p_{g,t}^{\text{total}} = P_g^{\min} + \textstyle\sum_s \delta_{g,s,t} && \forall g, t \\
    & && 0 \leq \delta_{g,s,t} \leq \bar{\delta}_g && \forall g, s, t \\
    & && p_{g,t}^{\text{total}} \geq P_g^{\min} && \forall g \in \mathcal{G}_{\text{must}}, t \tag{P4} \\
    & && -\text{RD}_g \leq p_{g,t}^{\text{total}} - p_{g,t-1}^{\text{total}} \leq \text{RU}_g && \forall g, t \geq 2 \tag{P5}
\end{align}

\subsection{Post-revision ED (Piecewise Linear + RE Bidding)}

\begin{align}
    &\min_{\delta, r, \text{re\_net}} &&
        \sum_{t,g,s} \text{mc}_{g,s} \cdot \delta_{g,s,t}
        + \sum_{t,g} A_{g,t} \cdot p_{g,t}^{\text{total}}
        + \sum_{t,k} b_{k,t} \cdot r_{k,t}
        + \varepsilon \sum_t \text{re\_net}_t \tag{R1} \\
    &\text{s.t.} &&
        \sum_g p_{g,t}^{\text{total}} + \sum_k r_{k,t} + \text{re\_net}_t = D_t && \forall t \tag{R2} \\
    & && P_{k,t}^{\min,\text{RE}} \leq r_{k,t} \leq \bar{R}_{k,t} && \forall k, t \tag{R3} \\
    & && 0 \leq \text{re\_net}_t \leq \text{RE}_t^{\text{nonbid}} && \forall t \\
    & && \text{(P4), (P5) 동일} \notag
\end{align}

\subsection{Price Adder 보정}

\begin{align}
    &\text{갱신:} && A_{g,t} \leftarrow A_{g,t} + \eta \cdot \frac{e_t}{n_{\text{marg},t}} && g \in \mathcal{M}_t \tag{C1} \\
    &\text{수축:} && A \leftarrow (1 - \lambda\eta) \cdot A \tag{C2} \\
    &\text{제한:} && A_{g,t} \in [-\bar{A}_g, \bar{A}_g] \tag{C3}
\end{align}

\subsection{Beta 몬테카를로}

\begin{align}
    &u_j \sim \text{Beta}(\alpha_j, \beta_j), \qquad
    \kappa \sim \max(0, \mathcal{N}(1, \sigma^2)) \tag{MC1} \\
    &\bar{u}_l = \frac{\sum_j \phi_j w_{j,l} u_j}{\sum_j \phi_j w_{j,l}} \tag{MC2} \\
    &b_l = \text{clamp}\!\big(\kappa(\bar{u}_l - 1) |\text{BidFloor}|, \text{BidFloor}, 0\big) \tag{MC3}
\end{align}


% ============================================================
% 10. 대표일 선정 방법론
% ============================================================
\section{대표일 선정 방법론}\label{sec:representative_days}

\subsection{일별 프로파일 특성}

각 날짜에 대해 다음 5가지 특성을 산출한다:
\begin{align}
    &\text{max\_demand}(d) = \max_{t \in d} D_t \\
    &\text{mean\_smp}(d) = \frac{1}{24}\sum_{t \in d} \text{SMP}_t \\
    &\text{solar\_share}(d) = \frac{\sum_{t \in d} \text{PV}_t}{\sum_{t \in d} D_t} \\
    &\text{wind\_share}(d) = \frac{\sum_{t \in d} \text{W}_t}{\sum_{t \in d} D_t} \\
    &\text{evening\_ramp}(d) = D_{21} - D_{15} \qquad [\text{MW}]
\end{align}

\subsection{계절별 3일 선정}

각 계절 $s$에 대해 3일을 선정한다:
\begin{enumerate}[noitemsep]
    \item \textbf{RE 고비율일}: 수요 중위수 이하 일자 중 $\text{solar\_share} + \text{wind\_share}$가 최대인 일자
    \item \textbf{SMP 평균일}: 나머지 일자 중 $|\text{mean\_smp} - \text{median}(\text{mean\_smp}_s)|$이 최소인 일자
    \item \textbf{피크 수요일}: $\text{max\_demand}$가 최대인 일자
\end{enumerate}
결과적으로 4계절 $\times$ 3일 $= 12$일이 test 대표일로 선정된다.


% ============================================================
% 11. 모형 한계 및 가정
% ============================================================
\section{모형 한계 및 가정}\label{sec:limitations}

\subsection{주요 가정}

\begin{enumerate}[noitemsep]
    \item \textbf{UC 미모형}: 기동정지(Unit Commitment)를 직접 모형화하지 않고
          Price Adder로 근사한다. 기동비, 무부하비, 최소가동시간의 시계열 효과는
          부분적으로만 반영된다.

    \item \textbf{단일 노드}: 송전 제약을 고려하지 않는 단일 노드(copper plate) 모형이다.
          송전혼잡(transmission congestion)에 의한 지역별 SMP 차이는 반영되지 않는다.

    \item \textbf{송전손실계수}: TLF(Transmission Loss Factor)를 1.0으로 가정한다.

    \item \textbf{RE Potential $\equiv$ 발전량}: 2024년 데이터에서 peak CF $< 1$이므로
          발전량을 physical availability로 직접 사용한다.

    \item \textbf{VOM 미포함}: 한국 CBP 시장에서 변동운영비(VOM)는 별도 정산되므로
          발전비용에 포함하지 않는다.

    \item \textbf{확정적 수요}: 수요는 실측 데이터를 사용하며 불확실성을 고려하지 않는다.
\end{enumerate}

\subsection{향후 확장 방향}

\begin{itemize}[noitemsep]
    \item UC(Unit Commitment) 통합 모형으로의 확장
    \item 다노드 모형 및 송전 제약 반영
    \item 에너지저장장치(ESS) 입찰 참여 모형화
    \item 수요 불확실성의 확률적 모형화
    \item REC 가격의 동태적 모형화
\end{itemize}


% ============================================================
% 부록
% ============================================================
\appendix

\section{Piecewise Linear 근사의 정확성}\label{app:piecewise}

2차 비용함수 $C(P) = aP^2 + bP + c$를 $S$개 구간으로 분할할 때,
각 구간의 근사 오차 한계를 분석한다.

구간 $s$의 폭 $\bar{\delta} = (P^{\max} - P^{\min})/S$에서
최대 오차는 구간 양 끝점과 중점의 MC 차이에 비례한다:
\begin{equation}
    \max_s |MC(p) - \text{mc}_s| \leq a \cdot \bar{\delta} \times 1000 \qquad [\text{원/MWh}]
\end{equation}
$S = 4$일 때 출력 범위 500 MW, $a = 0.001$인 경우:
$\bar{\delta} = 125$ MW, 최대 오차 $= 0.001 \times 125 \times 1000 = 125$ 원/MWh.

구간 수 $S$를 증가시키면 오차는 $O(1/S)$로 감소하지만, LP 크기가 비례 증가한다.
$S=4$는 정확도와 계산 효율의 균형점으로 선택되었다.


\section{LP Dual과 SMP의 관계}\label{app:dual_smp}

선형계획법(LP)의 강쌍대정리(Strong Duality Theorem)에 의해,
최적 원시해(primal)와 쌍대해(dual)의 목적값이 일치한다.

수급균형 제약 $\sum_g p_{g,t} = D_t$에 대응하는 쌍대변수 $\lambda_t$는:
\begin{equation}
    \lambda_t = \frac{\partial z^*}{\partial D_t}
\end{equation}
즉, 수요 $D_t$가 1 MW 증가할 때 최소 총비용 $z^*$의 증가분을 나타내며,
이는 경제학적으로 한계가격(marginal price)과 일치한다.

LP에서 한계발전기(marginal generator)가 유일하게 결정되면
$\lambda_t = c_{g^*}$ (한계발전기의 비용)이 된다.
복수의 한계발전기가 존재하면 LP 솔버에 의해 하나의 dual value가 선택된다.

Post-ED에서 $\text{re\_net}_t$ 변수의 비용 $\varepsilon_{\text{nonbid}}$는
$\text{re\_net}_t$가 상한에 도달하지 않는 한 dual에 영향을 주지 않는다.
상한에 도달하면(출력제한 발생) dual에 $\varepsilon_{\text{nonbid}}$가 반영되나,
이는 출력제한 시의 실제 비용 신호로 해석할 수 있다.


\section{원전 호기명-발전기 ID 매핑 테이블}\label{app:nuclear_map}

\begin{longtable}{ll}
\caption{원전 호기명 $\to$ Generator ID 매핑}\label{tab:nuclear_map} \\
\toprule
\textbf{호기명} & \textbf{Generator ID} \\
\midrule
\endfirsthead
\multicolumn{2}{c}{\tablename\ \thetable{} -- (이어서)} \\
\toprule
\textbf{호기명} & \textbf{Generator ID} \\
\midrule
\endhead
\bottomrule
\endfoot
새울1호기   & Nuclear\_001 \\
고리3호기   & Nuclear\_002 \\
고리4호기   & Nuclear\_003 \\
신고리1호기 & Nuclear\_004 \\
신고리2호기 & Nuclear\_005 \\
새울2호기   & Nuclear\_006 \\
새울3호기   & Nuclear\_007 \\
한빛1호기   & Nuclear\_008 \\
한빛2호기   & Nuclear\_009 \\
한빛3호기   & Nuclear\_010 \\
한빛4호기   & Nuclear\_011 \\
한빛5호기   & Nuclear\_012 \\
한빛6호기   & Nuclear\_013 \\
신한울1호기 & Nuclear\_014 \\
신한울2호기 & Nuclear\_015 \\
신월성2호기 & Nuclear\_016 \\
신월성1호기 & Nuclear\_017 \\
월성2호기   & Nuclear\_018 \\
월성3호기   & Nuclear\_019 \\
월성4호기   & Nuclear\_020 \\
한울1호기   & Nuclear\_021 \\
한울2호기   & Nuclear\_022 \\
한울3호기   & Nuclear\_023 \\
한울4호기   & Nuclear\_024 \\
한울5호기   & Nuclear\_025 \\
한울6호기   & Nuclear\_025 \\
\end{longtable}

\noindent 한울5호기와 한울6호기는 통합 모델(fallback)로 처리되어 동일 ID를 사용한다.
MATPOWER KPG193 데이터의 bus 번호 구분:
bus 82 = 고리/신고리/새울, bus 124 = 한빛,
bus 166 = 월성/신월성, bus 175 = 한울/신한울.


\section{기본 파라미터 값}\label{app:default_params}

\begin{table}[H]
\centering
\caption{주요 기본 파라미터}\label{tab:default_params}
\small
\begin{tabular}{llll}
\toprule
\textbf{파라미터} & \textbf{기호} & \textbf{기본값} & \textbf{비고} \\
\midrule
PW 세그먼트 수 & $S$ & 4 & 구간별 선형 근사 \\
입찰참여율 (PV) & $\rho_{\text{PV}}$ & 0.3 & \\
입찰참여율 (Wind) & $\rho_{\text{W}}$ & 0.3 & \\
하한가 계수 & $\beta$ & 2.0 & \\
REC 가격 & & 80 원/kWh & \\
RE Pmin 비율 & $\alpha$ & 0.1 & \\
비입찰 RE 비용 & $\varepsilon_{\text{nonbid}}$ & 100 원/MWh & \\
학습률 & $\eta$ & 0.3 (단일일) / 0.2 (다일) & \\
L2 shrinkage & $\lambda\eta$ & 0.05 & \\
단일일 최대 반복 & $N_{\max}$ & 20 & \\
다일 에폭 수 & & 10 & \\
목표 MAE & & 5,000 (단일일) / 4,000 (다일) & 원/MWh \\
Common shock $\sigma$ & $\sigma$ & 0.10 & \\
MC 표본 수 & $N$ & 200 (기본) / 100 (침투도) & \\
Buffer 일수 & & $\pm 3$ & Train/Test 분리 \\
계절당 Train 일수 & & 25 & \\
계절당 Test 일수 & & 3 & \\
블록 가중치 & $\bm{w}$ & (0.4, 0.3, 0.3) & Low/Mid/High \\
\bottomrule
\end{tabular}
\end{table}


\section{Version History}\label{app:changelog}

\begin{description}[style=nextline]
    \item[v2 (개별 발전기 122기 + real-data)]
    주요 변경사항:
    \begin{itemize}[noitemsep]
        \item 개별 발전기 122기 (클러스터링 없음), MATPOWER KPG193 기반
        \item heat\_rate/vom 제거: 한국 CBP 시장에서 gencost에 열소비율 내포
        \item Potential 재구성 (가정: 2024 발전량 $\equiv$ potential, max CF $< 1$)
        \item BidderType mixture (aggressive/moderate/conservative/PPA\_locked)
        \item Case\_A\_zero $\equiv$ $\rho=0$ baseline (\texttt{bidding\_active=false}), $\varepsilon_{\text{nonbid}}=100$
        \item RE Pmin = $\min(\alpha \cdot \text{installed\_mw}, \text{avail})$
        \item Beta$(\alpha,\beta)$ mixture + common shock 몬테카를로
        \item 활성 marginal 전체 갱신 + $1/n_{\text{marg}}$ 정규화
        \item Tikhonov L2 shrinkage ($\lambda\eta = 0.05$/iter)
        \item Curtailment-free calibration (RE must-take, dual purity)
        \item Train/test/buffer split (계절 $\pm$3일 buffer) + 3D adder $G \times 24 \times S$
        \item Sanity assert: $|\text{SMP}_{\text{post},A} - \text{SMP}_{\text{pre}}| < 1$
    \end{itemize}
\end{description}


\end{document}
